\usepackage{xcolor}
\usepackage{listings}

% 1. IMPOSTAZIONE FONT SANS-SERIF (Helvetica/Arial-like)
\usepackage{helvet}
\renewcommand{\familydefault}{\sfdefault}

% 2. COLORAZIONE TITOLI (SECTIONS IN BLU)
\usepackage{sectsty}
\definecolor{sectionblue}{RGB}{0,102,204} % Definisci la tonalità di blu desiderata
\definecolor{sectiongreen}{RGB}{0,110,0} % Definisci la tonalità di blu desiderata
\sectionfont{\color{sectionblue}}          % Applica il colore a \section
\subsectionfont{\color{sectiongreen}}       % (Opzionale) Applica anche a \subsection
\sectionfont{\color{sectionblue}}          % Applica il colore a \section

\definecolor{backcolour}{RGB}{248,248,248}
\definecolor{codegreen}{RGB}{0,110,0}
\definecolor{codegray}{RGB}{100,100,100}
\definecolor{codepurple}{RGB}{125,0,125}
\definecolor{codesyntax}{RGB}{0,60,160}

\lstdefinestyle{cstyle}{
  language=C,
  basicstyle=\ttfamily\footnotesize,
  backgroundcolor=\color{backcolour},
  commentstyle=\color{codegreen},
  keywordstyle=\color{codesyntax},
  numberstyle=\tiny\color{codegray},
  stringstyle=\color{codepurple},
  breakatwhitespace=false,
  breaklines=true,
  captionpos=t,
  keepspaces=true,
  numbers=left,
  numbersep=5pt,
  showspaces=false,
  showstringspaces=false,
  showtabs=false,
  tabsize=2
}

% Applica lo stile ai blocchi di codice generati da Pandoc.
\lstset{style=cstyle}